% -------------------------------------------------------
\section{Listados de Código}\label{sec:codigo}

Texforge resalta el código fuente en Rust durante la compilación: cada
bloque \texttt{code} se reescribe sobre la copia de compilación y solo
necesita \texttt{color.sty}, sin depender de \texttt{listings} ni de
\texttt{minted}. El lenguaje se indica con \texttt{lang} y la numeración de
líneas con \texttt{numbers}. El tema tipográfico ---\texttt{one-light} en
este caso--- se elige una sola vez para todo el documento en
\texttt{project.toml}, y el \emph{estilo} de color se puede cambiar por
bloque: hay cuatro (\texttt{light}, \texttt{light-mono}, \texttt{dark} y
\texttt{dark-mono}) y se pueden mezclar en un mismo documento.

\subsection{Python}

Con \texttt{numbers=true} se añade una guía de líneas dentro del marco:

\begin{code}[lang=python, numbers=true]
import numpy as np

def integrar_gauss(n_puntos=1000):
    """Aproxima la integral de Gauss por cuadratura."""
    x = np.linspace(-10, 10, n_puntos)
    dx = x[1] - x[0]
    y = np.exp(-x**2)
    return np.sum(y) * dx

if __name__ == "__main__":
    print(f"Integral = {integrar_gauss():.6f}")
\end{code}

\subsection{Rust}

El mismo mecanismo resalta Rust, aquí con numeración: el
Listado~\ref{lst:fib} muestra la definición clásica de Fibonacci.

\begin{code}[lang=rust, numbers=true, caption={Fibonacci en Rust}, label={lst:fib}]
fn fibonacci(n: u64) -> u64 {
    match n {
        0 => 0,
        1 => 1,
        _ => fibonacci(n - 1) + fibonacci(n - 2),
    }
}

fn main() {
    println!("fib(10) = {}", fibonacci(10));
}
\end{code}

\subsection{Shell}

La tabla \texttt{[highlight.by\_lang]} de \texttt{project.toml} fija
\texttt{bash = "dark"}, así que este bloque usa el estilo \texttt{dark} sin
tener que pedirlo en su propio encabezado --- igual que una terminal. Estos
caracteres son especiales en LaTeX y se imprimen tal cual, porque cada
carácter se escapa antes de llegar al motor:

\begin{code}[lang=bash]
#!/usr/bin/env bash
set -euo pipefail

echo '$ % & # _ ^ ~ { } \'
\end{code}

\subsection{Estilos monocromos}

Un bloque puede pedir su propio estilo con la opción \texttt{style}. El
variante \texttt{light-mono} quita todo el color: las palabras clave van en
negrita y los comentarios en cursiva, que es lo que funciona en una
impresión en blanco y negro.

\begin{code}[lang=python, style=light-mono]
def sumar(numeros):  # solo grises: apto para B/N
    return sum(numeros)
\end{code}

\subsection{TOML}

La configuración del propio proyecto se muestra con el mismo entorno:

\begin{code}[lang=toml]
[highlight]
theme = "one-light"
style = "light"

[highlight.by_lang]
bash = "dark"
\end{code}

\listoflistings
